%% Die Sprachdatei von helsinki für de-DE (#377): die Wörter des Designs und
%% die Schreibweise der Zahlen, soweit sie an der Sprache hängt. Der Renderer
%% legt sie als sprache.tex neben die Vorlage, die sie lädt.
%%
%% Jede Sprachdatei eines Designs definiert dieselben Makros, mit derselben
%% Stelligkeit – das prüft ein Test. Welches Wort ein Design an einer Stelle
%% setzt, entscheidet das Design; deshalb eine Datei je Design und Sprache
%% und keine für alle Designs. Ein Satz mit fachlicher Aussage steht nicht
%% hier, sondern kommt aus Kotlin (ADR-0017). Kein Platzhalter: Die Werte der
%% Rechnung setzt die Vorlage.
%%
%% Die Beschriftungen heißen \wort... nach dem deutschen Wort, damit sie sich
%% von den Makros unterscheiden, die der Renderer aufruft: \wortbetrag ist
%% der Spaltenkopf, \betrag die Zahl.

%% ---------------------------------------------------------------- Zahlen
%% Dezimalkomma und Punkt als Tausendertrenner (ADR-0021). Ab wie vielen
%% Stellen gruppiert wird, wählt die Vorlage.
\sisetup{
  output-decimal-marker = {,},
  group-separator       = {.},
}
%% Das Währungszeichen und das Prozentzeichen hängen mit festem Leerzeichen
%% an der Zahl. Einzelpreise nehmen \betrag.
%%
%% Welches Zeichen, sagt die Währung der Rechnung (#409): Die Vorlage setzt
%% \waehrung auf ihren Code, und \waehrungszeichen schlägt damit ein Zeichen
%% je Code nach. Eines für jeden Code aus Currency, das prüft Template.of.
%% $ und nicht US$ oder USD: das Zeichen, das der Empfänger auf seiner
%% Rechnung erwartet; welcher Dollar gemeint ist, sagt das XML (BT-5).
\newcommand{\zeichenEUR}{€}
\newcommand{\zeichenUSD}{\$}
\newcommand{\zeichenGBP}{£}
\newcommand{\waehrungszeichen}{\csname zeichen\waehrung\endcsname}
\newcommand{\betrag}[1]{\num{#1}~\waehrungszeichen}
\newcommand{\prozent}[1]{\num{#1}~\%}

%% ------------------------------------------------------------- Seiten
\newcommand{\wortseitevon}[2]{Seite #1 von #2}
\newcommand{\wortfortsetzung}{Fortsetzung auf der nächsten Seite}

%% ------------------------------------------------------ Kopf und Anschrift
%% "Geschäftsführung" statt "Geschäftsführer": Die Bezeichnung steht neben
%% einem Namen, dessen Geschlecht die Stammdaten nicht kennen.
\newcommand{\wortgeschaeftsfuehrung}{Geschäftsführung}
\newcommand{\wortustid}{USt-IdNr.}
\newcommand{\worttelefon}{Tel.}
\newcommand{\wortkundennr}{Kundennr.:}
\newcommand{\wortauftragsnr}{Unsere Auftragsnr.:}
\newcommand{\wortvereinnahmtam}{Abschlag vereinnahmt am}
\newcommand{\wortreferenz}{Ihre Referenz:}
\newcommand{\wortzahlbarbis}{Zahlbar ohne Abzug bis}

%% ----------------------------------------------------- Positionen, Summen
\newcommand{\wortbeschreibung}{Beschreibung}
\newcommand{\wortmenge}{Menge}
\newcommand{\worteinzelpreis}{Einzelpreis}
\newcommand{\wortrabatt}{Rabatt}
\newcommand{\wortbetrag}{Betrag}
\newcommand{\wortnettobetrag}{Nettobetrag}
%% Satz und, bei mehreren Steuergruppen, ihre Grundlage – sonst ist sie leer.
\newcommand{\wortsteuer}[2]{zzgl. #1 USt.\ifblank{#2}{}{ auf #2}}
\newcommand{\wortgesamtbetrag}{Gesamtbetrag}

%% --------------------------------------------------------- Schluss und Fuß
\newcommand{\wortgruss}{Mit freundlichen Grüßen}
\newcommand{\wortbankverbindung}{Bankverbindung}
\newcommand{\wortkontoinhaber}{Kontoinhaber}
\newcommand{\wortinstitut}{Institut}
\newcommand{\wortregistergericht}{Registergericht}
